\chapter{La completación de un espacio métrico}\label{chap:09}
\section{Agregar límites sin alterar las distancias}
Cuando un espacio no es completo, podemos intentar añadir los límites que le faltan. Hay dos exigencias: las distancias entre los puntos originales deben conservarse y los puntos nuevos deben surgir como límites de puntos antiguos. La construcción siguiente satisface ambas y, además, es única salvo una identificación que conserva distancias.

\begin{definicion}
Una aplicación $i:(X,d)\to(Y,\rho)$ es una \emph{isometría} si $\rho(i(x),i(y))=d(x,y)$ para todos $x,y\in X$. Una \emph{completación} de $X$ es un espacio completo $\widehat X$ junto con una isometría $i:X\to\widehat X$ cuya imagen es densa.
\end{definicion}
\index{isometria@isometría}\index{completacion@completación}
Una isometría es inyectiva: imágenes iguales dan distancia cero entre los puntos originales. Es $1$-Lipschitz y, por tanto, uniformemente continua; su inversa sobre la imagen también es una isometría. La inclusión de un subespacio con la distancia heredada es un ejemplo. Si $X$ ya es completo, $(X,\id_X)$ es una completación suya.

\section{Construcción mediante sucesiones de Cauchy}
\begin{teorema}\label{ch09:existencia}
Todo espacio métrico tiene una completación.
\end{teorema}
\begin{proof}
El espacio vacío es completo y se completa con él mismo. Supongamos $X\ne\varnothing$ y sea $\mathcal C(X)$ el conjunto de todas las sucesiones de Cauchy de $X$. Este conjunto no es vacío porque contiene todas las sucesiones constantes.

\textbf{Paso 1: una distancia provisional.} Para $u=(u_n)$ y $v=(v_n)$ ponemos
\[
 D(u,v)=\lim_{n\to\infty}d(u_n,v_n).
\]
Debemos demostrar que el límite existe. La desigualdad triangular, aplicada en los dos sentidos, produce
\[
 |d(u_n,v_n)-d(u_m,v_m)|\leq d(u_n,u_m)+d(v_n,v_m).
\]
Las dos sucesiones del lado derecho son de Cauchy; así, $(d(u_n,v_n))$ es una sucesión real de Cauchy. Su límite existe por la completitud de $\R$. Además, $D\geq0$, $D(u,u)=0$ y $D(u,v)=D(v,u)$. Si $w=(w_n)$, de $d(u_n,v_n)\leq d(u_n,w_n)+d(w_n,v_n)$ y la preservación de desigualdades al pasar al límite resulta $D(u,v)\leq D(u,w)+D(w,v)$. Por ello $D$ es una pseudométrica; puede anularse sobre sucesiones diferentes, por ejemplo sobre dos sucesiones que converjan al mismo punto.

\textbf{Paso 2: identificar las representaciones del mismo punto.} Definimos $u\sim v$ cuando $D(u,v)=0$. La reflexividad y la simetría provienen de las de $D$; la transitividad se sigue de $D(u,w)\leq D(u,v)+D(v,w)$. Sea
\[
 \widehat X=\mathcal C(X)/\!\sim,
 \qquad \widehat d([u],[v])=D(u,v).
\]
Esta fórmula no depende de representantes: si $u\sim u'$ y $v\sim v'$, entonces
\[
 |D(u,v)-D(u',v')|\leq D(u,u')+D(v,v')=0.
\]
Las propiedades de pseudométrica descienden al cociente y, además, $\widehat d([u],[v])=0$ implica $[u]=[v]$. Tenemos una métrica verdadera.

\textbf{Paso 3: insertar $X$ y probar densidad.} Definimos $i(x)$ como la clase de la sucesión constante $(x,x,\ldots)$. Entonces $\widehat d(i(x),i(y))=\lim_nd(x,y)=d(x,y)$, por lo que $i$ es una isometría. Si $u=[(u_n)]$, se cumple
\[
 \widehat d(i(u_n),u)=\lim_{k\to\infty}d(u_n,u_k)\longrightarrow0
 \quad(n\to\infty).
\]
En efecto, dado $\eps>0$, la condición de Cauchy permite que $d(u_n,u_k)<\eps/2$ para $n,k\geq N$; fijando $n\geq N$ y pasando $k$ al límite, queda una distancia no mayor que $\eps/2<\eps$. Esto prueba que $i(X)$ es denso.

\textbf{Paso 4: completar la construcción.} Sea $(U_n)$ una sucesión de Cauchy en $\widehat X$. Por la densidad ya probada, para cada $n$ elegimos $z_n\in X$ con
\[
 \widehat d(i(z_n),U_n)<2^{-n}.
\]
Estas elecciones se hacen una sola vez, con tolerancias predeterminadas. No dependen del $\eps$ que aparecerá después al comprobar la convergencia. Para $n,m\in\N$,
\[
 d(z_n,z_m)\leq 2^{-n}+\widehat d(U_n,U_m)+2^{-m}.
\]
Dado $\eps>0$, elegimos $N$ de modo que los dos términos geométricos sumen menos que $\eps/2$ para $n,m\geq N$ y que $\widehat d(U_n,U_m)<\eps/2$. Así $(z_n)$ es de Cauchy en $X$. Sea $Z=[(z_n)]\in\widehat X$. El paso 3 aplicado a esta sucesión da $i(z_n)\to Z$. Finalmente,
\[
 \widehat d(U_n,Z)\leq\widehat d(U_n,i(z_n))+\widehat d(i(z_n),Z)
 <2^{-n}+\widehat d(i(z_n),Z)\longrightarrow0.
\]
Toda sucesión de Cauchy de $\widehat X$ converge allí; por tanto es completo.
\end{proof}

\begin{observacion}[La elección diagonal y sus cuantificadores]
Si cada $U_n$ se representa por una sucesión $(u_{n,k})_k$, el punto $z_n$ del paso 4 puede tomarse como $u_{n,k_n}$ con $k_n$ suficientemente grande. Se puede imponer además $k_n>k_{n-1}$, porque la aproximación vale para todos los índices de una cola. Lo esencial es escoger $\widehat d(i(u_{n,k_n}),U_n)<2^{-n}$, y no un error fijo $\eps$ antes de definir la diagonal: una sucesión candidata al límite debe estar definida independientemente de la tolerancia con la que más tarde se compruebe que converge.
\end{observacion}

\section{Dos modelos y una precaución sobre las normas}
\begin{ejemplo}[Aproximación de funciones discontinuas]
En $C([0,1])$, la fórmula $\norm{f}_1=\int_0^1|f(t)|\,dt$ define una norma. Si una función continua no es cero, su valor absoluto es mayor que una constante positiva en algún intervalo de longitud positiva, y su integral es positiva. El valor absoluto es indispensable: $\int f$ por sí sola no es una norma.

Este espacio normado no es completo. Para $n\geq3$, definamos $f_n$ igual a cero en $[0,1/2]$, lineal de cero a uno en $[1/2,1/2+1/n]$ e igual a uno en el intervalo restante. Si $h$ es cero en $[0,1/2]$ y uno en $(1/2,1]$, entonces
\[
 \int_0^1|f_n-h|\leq1/n,
 \quad\text{y por tanto}\quad
 \norm{f_n-f_m}_1\leq1/n+1/m.
\]
Si existiera un límite continuo $f$ en esta norma, la misma desigualdad triangular daría $\int|f-h|=0$. En cualquier intervalo contenido en $(0,1/2)$, la continuidad obligaría a $f=0$: un valor no nulo produciría una integral positiva en una vecindad. De igual forma, $f=1$ en $(1/2,1)$. Esto contradice la continuidad en $1/2$.

El espacio de clases de funciones Lebesgue-integrables $L^1([0,1])$ proporcionará una realización de esta completación en el volumen de medida e integración. Aquí no se presupone esa teoría: la completación por sucesiones ya está construida. Para un compacto arbitrario $K\subset\R^n$, la integral puede ser sólo una seminorma sobre $C(K)$; por ejemplo, si $K=\{0\}$, todas las integrales respecto del volumen son cero. Por ello la afirmación sobre $C(K)$ requiere hipótesis sobre la medida y su soporte, o pasar primero al cociente por las funciones de seminorma cero.
\end{ejemplo}

\begin{ejemplo}[Los números p-ádicos]\label{ch09:qp}
La completación de $(\Q,d_p)$ se denota por $\Q_p$. Para una clase $r=[(r_n)]$, definimos $|r|_p=\widehat d(r,0)=\lim_n|r_n|_p$. La desigualdad ultramétrica pasa al límite, de modo que sigue siendo válida en la completación. Además de ser un espacio completo, conserva las operaciones de cuerpo. Veamos por qué. Si $(r_n)$ y $(s_n)$ son de Cauchy, sus valores absolutos p-ádicos están acotados, y
\[
 |r_ns_n-r_ms_m|_p\leq
 \max\{|r_n|_p|s_n-s_m|_p,\,|s_m|_p|r_n-r_m|_p\}.
\]
Esto prueba que los productos forman una sucesión de Cauchy; para la suma, $|(r_n+s_n)-(r_m+s_m)|_p$ está acotado por el máximo de las dos diferencias correspondientes. Las mismas estimaciones muestran que cambiar representantes por sucesiones equivalentes no cambia la clase de su suma o producto. Las identidades de cuerpo se conservan porque valen término a término. También $|rs|_p=|r|_p|s|_p$, pues la identidad racional $|r_ns_n|_p=|r_n|_p|s_n|_p$ pasa al límite de números reales.

Falta justificar la inversa. Para una clase $r\ne0$, su distancia a cero es un número $c>0$. Si $(r_n)$ la representa, entonces $|r_n|_p\to c$, de modo que $|r_n|_p\geq c/2$ en una cola. Allí
\[
 |r_n^{-1}-r_m^{-1}|_p
 =\frac{|r_n-r_m|_p}{|r_n|_p|r_m|_p}
 \leq \frac{4}{c^2}|r_n-r_m|_p.
\]
La sucesión de inversas es de Cauchy y su clase multiplica a $r$ para dar $1$. Los términos iniciales se pueden reemplazar sin alterar la clase. Así $\Q_p$ es un cuerpo.

Todo elemento tiene una expansión
\[
 r=\sum_{k\geq k_0}a_kp^k,
 \qquad k_0\in\Z,\quad a_k\in\{0,\ldots,p-1\},
\]
entendida como límite p-ádico de sus sumas parciales. Damos la razón de la existencia. Primero supongamos $|r|_p\leq1$ y tomemos racionales $r_n\to r$. Para cada $m\geq1$, elegimos uno con $|r-r_n|_p<p^{-m}$ y $|r_n|_p\leq1$. Podemos escribirlo $r_n=a/b$ con $p\nmid b$. Como $b$ tiene inversa módulo $p^m$, existe un único entero $c_m\in\{0,\ldots,p^m-1\}$ con $bc_m\equiv a\pmod{p^m}$. Por tanto $|r-c_m|_p\leq p^{-m}$. La unicidad de estos residuos, y la desigualdad ultramétrica, implican $c_{m+1}\equiv c_m\pmod{p^m}$; en consecuencia $c_{m+1}=c_m+a_mp^m$ para un dígito $a_m$. Partiendo de $c_1=a_0$, obtenemos la expansión y la cota $|r-c_m|_p\leq p^{-m}\to0$.

Para un $r$ arbitrario, su norma es finita y podemos multiplicarlo por una potencia positiva de $p$ hasta que su norma sea a lo sumo $1$; deshacer esa multiplicación produce el índice inicial $k_0$ posiblemente negativo. La unicidad de los dígitos se demuestra como en el ejemplo~\ref{ch08:padica-incompleta}: el primer dígito distinto determina la valoración de la diferencia y no puede ser cancelado por términos de mayor potencia. Este argumento muestra también cómo la completación agrega puntos que la aritmética racional no contenía.
\end{ejemplo}

\section{Ejercicios resueltos}
\begin{ejercicio}[Lectura 9, ejercicio 1]\label{ch09:unicidad}
Sean $(Y,i)$ y $(Y',i')$ dos completaciones de $X$. Demuestre que existe una isometría sobreyectiva $T:Y\to Y'$ tal que $T\circ i=i'$.
\end{ejercicio}
\begin{solucion}
Si $X=\varnothing$, la densidad fuerza $Y=Y'=\varnothing$, y la afirmación es inmediata. En otro caso, para $y\in Y$ elegimos $x_n\in X$ con $i(x_n)\to y$. Como $i$ es isometría, $(x_n)$ es de Cauchy en $X$, y $(i'(x_n))$ lo es en el espacio completo $Y'$. Definimos
\[
 T(y)=\lim_{n\to\infty}i'(x_n).
\]
Si otra sucesión $z_n$ satisface $i(z_n)\to y$, entonces
\[
 d(x_n,z_n)=\rho(i(x_n),i(z_n))
 \leq\rho(i(x_n),y)+\rho(y,i(z_n))\longrightarrow0.
\]
La distancia entre $i'(x_n)$ e $i'(z_n)$ es la misma; sus límites coinciden y $T$ está bien definida. Para $y,z\in Y$, elegimos aproximaciones $i(x_n)\to y$, $i(z_n)\to z$ y obtenemos
\[
 \rho'(T(y),T(z))=\lim_nd(x_n,z_n)=\rho(y,z).
\]
Aquí la continuidad de la distancia se justifica por $|\rho(a,b)-\rho(c,d)|\leq\rho(a,c)+\rho(b,d)$. Por tanto $T$ es una isometría. La sucesión constante $x_n=x$ muestra $T(i(x))=i'(x)$.

Para probar sobreyectividad, sea $y'\in Y'$. Elegimos $x_n\in X$ con $i'(x_n)\to y'$. Entonces $(i(x_n))$ es de Cauchy y converge en $Y$ a cierto $y$. Por la definición anterior, $T(y)=y'$. Finalmente, la isometría que cumple la identidad requerida es única: es continua y queda determinada en la imagen densa $i(X)$; si $S$ es otra, para $i(x_n)\to y$ tenemos $S(y)=\lim_ni'(x_n)=T(y)$.
\end{solucion}

\begin{corolario}[Propiedad de extensión]\label{ch09:extension}
Toda aplicación uniformemente continua $f:X\to Z$, con $Z$ completo, se extiende de manera única a una aplicación uniformemente continua $\widehat f:\widehat X\to Z$ tal que $\widehat f\circ i=f$.
\end{corolario}
\begin{proof}
La aplicación $f\circ i^{-1}:i(X)\to Z$ es uniformemente continua porque $i^{-1}$ conserva distancias. Aplicamos la construcción y las estimaciones del ejercicio~\ref{ch07:ej-extension} al conjunto denso $i(X)\subset\widehat X$. Su clausura es todo $\widehat X$, por lo que obtenemos exactamente la extensión solicitada y su unicidad.
\end{proof}

\section{La idea estructural}
Un punto nuevo es una clase de aproximaciones que terminan siendo indistinguibles en distancia. La densidad evita añadir puntos ajenos a esa idea, y la unicidad prueba que el resultado no depende de la realización concreta escogida. La completitud convierte después construcciones infinitas en puntos efectivos del espacio.
